\documentclass{beamer}
\usepackage{graphicx}
\usepackage{amsmath, amssymb}
\usepackage{qcircuit} % For quantum circuits
\usepackage{tikz}
\usetikzlibrary{quantikz} % Quantum circuit diagrams

\title{Quantum Phase Estimation}
\author{Morten HJ}
\date{Modified slides from 2009}

\begin{document}

\begin{frame}
    \titlepage
\end{frame}

% Slide: Introduction
\begin{frame}{Introduction to Quantum Phase Estimation (QPE)}
    \begin{itemize}
        \item QPE is a fundamental quantum algorithm used to estimate the phase $\phi$ in eigenvalue equations of unitary operators.
        \item Given a unitary operator $U$ and an eigenvector $| \psi \rangle$ such that:
        \[
        U | \psi \rangle = e^{2\pi i \phi} | \psi \rangle
        \]
        the goal of QPE is to estimate $\phi$.
        \item Applications:
        \begin{itemize}
            \item Shor's factoring algorithm
            \item Hamiltonian simulation
            \item Quantum chemistry
        \end{itemize}
    \end{itemize}
\end{frame}

% Slide: Quantum Circuit
\begin{frame}{Quantum Circuit for QPE}
    \centering
    \begin{quantikz}
        \lstick[wires=1]{$n$ qubits} & \gate{H} & \qwbundle{n} & \ctrl{1} & \qw & \qw & \gate{\text{QFT}^{\dagger}} & \meter{} \\
        \lstick[wires=1]{1 qubit} & \qw & \qw & \gate{U^{2^j}} & \qw & \qw & \qw & \qw
    \end{quantikz}
    \begin{itemize}
        \item The first register contains $n$ qubits initialized to $|0\rangle$.
        \item Hadamard gates create a superposition state.
        \item Controlled-$U^{2^j}$ operations apply phase shifts.
        \item The inverse Quantum Fourier Transform (QFT) extracts phase information.
    \end{itemize}
\end{frame}

% Slide: State Evolution
\begin{frame}{State Evolution in QPE}
    \begin{itemize}
        \item Initial state:
        \[
        | \psi_0 \rangle = | 0 \rangle^{\otimes n} \otimes | \psi \rangle
        \]
        \item After Hadamard transform:
        \[
        \frac{1}{\sqrt{2^n}} \sum_{k=0}^{2^n - 1} | k \rangle | \psi \rangle
        \]
        \item Applying controlled-$U$ gates:
        \[
        \frac{1}{\sqrt{2^n}} \sum_{k=0}^{2^n - 1} | k \rangle U^k | \psi \rangle
        \]
        Since $U | \psi \rangle = e^{2\pi i \phi} | \psi \rangle$, this becomes:
        \[
        \frac{1}{\sqrt{2^n}} \sum_{k=0}^{2^n - 1} e^{2\pi i k \phi} | k \rangle | \psi \rangle
        \]
    \end{itemize}
\end{frame}

% Slide: Inverse QFT
\begin{frame}{Inverse Quantum Fourier Transform (QFT)}
    \begin{itemize}
        \item The QFT maps computational basis states to phase-encoded states:
        \[
        | k \rangle \to \frac{1}{\sqrt{2^n}} \sum_{j=0}^{2^n - 1} e^{2\pi i k j / 2^n} | j \rangle
        \]
        \item Applying inverse QFT reveals the phase $\phi$ as a binary fraction.
    \end{itemize}
    \centering
    \begin{tikzpicture}
%        \node (qft) at (0,0) {\includegraphics[width=0.7\textwidth]{qft_diagram.png}};
    \end{tikzpicture}
\end{frame}

% Slide: Example Computation
\begin{frame}{Example: Phase of a Simple Unitary}
    \begin{itemize}
        \item Consider $U = \begin{bmatrix} 1 & 0 \\ 0 & e^{2\pi i / 8} \end{bmatrix}$, which has eigenvalues $e^{2\pi i \phi}$ with $\phi = 1/8$.
        \item Using 3 qubits for precision:
        \[
        \phi = 0.001_2 = \frac{1}{8}
        \]
        \item Measurement should give result $\phi = 0.001$ (binary).
    \end{itemize}
\end{frame}

% Slide: Applications of QPE
\begin{frame}{Applications of QPE}
    \begin{itemize}
        \item \textbf{Shor's Algorithm}: Finds the period of a modular function.
        \item \textbf{Hamiltonian Simulation}: Extracts energy eigenvalues for molecular systems.
        \item \textbf{Quantum Metrology}: Improves phase measurement precision.
    \end{itemize}
\end{frame}

% Slide: Error Analysis
\begin{frame}{Error Analysis and Precision}
    \begin{itemize}
        \item The precision of QPE depends on the number of qubits used.
        \item Probability of failure decreases exponentially with $n$ qubits.
        \[
        P_{\text{error}} \approx \frac{1}{2^{2n}}
        \]
        \item Noise and decoherence can impact accuracy, requiring error correction.
    \end{itemize}
\end{frame}

% Slide: Conclusion
\begin{frame}{Conclusion}
    \begin{itemize}
        \item QPE is a key quantum subroutine for extracting eigenphases.
        \item Involves controlled unitary operations and inverse QFT.
        \item Essential for algorithms like Shor's and Hamiltonian simulation.
    \end{itemize}
\end{frame}

\end{document}
